\subsection{Structural criteria and remaining sharpness questions}
\label{tht:program}

The positive and negative results suggest examining four distinct
features of a normalized Hamiltonian, in this order.

\begin{enumerate}
\item \textbf{Thermally occupied sectors.} Determine their restricted
partition functions and their weights at the requested error scale.
For a rare hard sector of weight $w_n$, a lower bound is multiplied by
$w_n$. A fixed-error separation requires nonvanishing weight; when the
requested error also decreases, compare the lower bound with a
polynomial in its inverse. A statement at constant error does not
classify every temperature schedule.
\item \textbf{Information stored in a sector.} A magnetization phase,
or finitely many species magnetizations, can be learned by coordinate
sampling. The bent label of a character code has a small quantum query
cost but needs an unbounded number of classical queries. Sector count,
covariance, entropy density, and a spectral gap of the Hamiltonian do
not decide this distinction by themselves.
\item \textbf{Conditional access.} Look for a product or other
simulable source whose output oracle is cheap to implement, and whose
source can be sampled conditionally on the observed input using few
input queries. Forward sampling alone supplies only the first half.
For an Ising tree, conditioning all spins factorizes the latent edge
variables, making the second half possible. Alternatively, establish
the reveal-progress criterion uniformly over pinnings.
\item \textbf{Geometry of interactions in the observed basis.}
Bounded interaction order, incident norm, loops, frustration, and
measurement basis must be recorded separately. A bounded single-spin
energy change does not bound the total effect of the other spins on a
site's conditional law. The global noisy-code obstruction exploits
exactly this difference. Cycles are not by themselves a negative
criterion: independent finite blocks and permutation-symmetric dense
models remain positive examples.
\end{enumerate}

For a classical Ising model use, for example,
\[
 H_n(x)=-\sum_{i<j}J_{ij}x_ix_j-\sum_i h_i x_i,
 \qquad \alpha(\beta)=\max_i\sum_{j\ne i}
                  \tanh(\beta|J_{ij}|).
\]
\textbf{Open conjecture (sufficient high-temperature regime).}
For every fixed $\kappa>0$, the family of such models with
$\alpha(\beta)\le1-\kappa$ satisfies full AA uniformly in dimension
and external fields, with polynomial dependence on $T,1/\epsilon$
for fixed $\kappa$. A stronger quantitative version asks for a bound
polynomial also in $1/\kappa$. These are conjectures here, not imported
consequences of uniqueness or mixing.

The later resolution proves an absolute subregion of this conjecture:
$\beta\max_i\sum_j|J_{ij}|\le1/8$ suffices, with no degree bound
and arbitrary fields. More generally,
Theorem~\ref{thm:resolution-dynamics-arbitrary-degree} uses the
weighted incident sum $\beta\max_i\sum_{A\ni i}(|A|-1)|J_A|$.
It does not establish the whole stated Dobrushin regime.

A subsequent derivation in
Section~\ref{resolution-dynamics:dobrushin-attempt} improves the
pairwise sufficient bound to $\beta\max_i\sum_j|J_{ij}|\le1/5$,
with the same $O(T^{18}\epsilon^{-28})$ query exponents. It separates
conditional forward branching, unconditional reverse branching, and
posterior rejection costs. That section also proves that the
single-update posterior support contraction used in this approach
can fail for every random-map representation even when the
Dobrushin coefficient is $2/5$. This is an obstruction to that proof
step, not a counterexample to the Dobrushin simulation conjecture.
The new derivation is separate from the earlier audited revision.

For measured quantum Gibbs laws, the bounded-locality and bounded
per-site term-norm hypothesis is now sufficient at an explicit small
constant: Corollary~\ref{cor:quantum-clusters-full-aa} proves the
result for noncommuting interactions and binary product POVMs. The
proof controls the effective classical potential, then supplies both
forward query access and local posterior access. The earlier thermal
preparation and correlation results alone were insufficient to supply
those steps.

The clearest proof target is therefore not ``mixing implies AA,'' but
one of two concrete bridges: derive polynomial reveal progress under
every pinning, or construct forward oracle access together with local
posterior access and a dimension-free quantum hybrid error bound.
Forests demonstrate the latter bridge even beyond a uniqueness regime.
Bounded-degree loopy Ising models now have a proved high-temperature
regime, and explicit signed pairwise families fail at low temperature.
A sharp boundary for those local counterexamples remains unproved.

The classification extension in
Section~\ref{sec:classification-tree-partitions} now supplies the
forward and posterior construction on a tree of bounded-size bags.
It proves all-finite-temperature simulation for cactus graphs,
fixed-width strips, and bounded-treewidth graphs with bounded degree.
The table in Section~\ref{sec:hamiltonian-classification} records
these classes alongside exact global-code thresholds and the
remaining sufficient-only regimes.

Finally, Section~\ref{sec:resolution-nonmonotone} disproves general
monotonicity: one extensive Hamiltonian has full AA exactly on
$[0,1)\cup(2,\infty]$. Thus the dossier establishes all-temperature
simulation, a finite positive threshold, simulation only at $\beta=0$
among fixed nonnegative inverse temperatures, and a disconnected
positive temperature set. No universal thermodynamic critical point
describes these possibilities.
